\section{Why a finite range can remove the apparent flag-variance penalty}
\label{sec:weak-dispersion-range}

The earlier weak-dispersion limit compares the fractional observation with
its canonical unit-range flag. It is not a lower bound for all quantum
encodings. I now give a uniform support argument explaining why a competent
single flag can recover the relevant small scale on a bounded numerical
input family. This argument is a checked consequence of exponential
remainders, convexity and range normalisation, not a priority claim for those
methods.

\begin{proposition}[Uniform log-gap range under centred perturbations]
Let $w_i>0$, $\sum_iw_i=1$, $H>0$ and $0\leq\delta_\varepsilon<\infty$, and let real centred perturbations
$h_i(\varepsilon,v)$ satisfy
$\sum_iw_i h_i(\varepsilon,v)=0$ and
$|h_i(\varepsilon,v)|\leq\delta_\varepsilon$ for every point $v$ in a
specified input support. Put
\[
 J_\varepsilon(v)=\log\sum_iw_i e^{h_i(\varepsilon,v)}.
\]
Then
\[
 0\leq J_\varepsilon(v)\leq
 {e^{\delta_\varepsilon}\over2}
       \sum_iw_i h_i(\varepsilon,v)^2
 \leq{e^{\delta_\varepsilon}\over2}\delta_\varepsilon^2
\]
uniformly on that support. If an increasing scalar conditional price $P$
has a finite global Lipschitz constant $L\geq0$, then its normalised controlled residual
$X_\varepsilon=[P(m+J_\varepsilon)-P(m)]/H$ satisfies
$0\leq X_\varepsilon\leq Le^{\delta_\varepsilon}\delta_\varepsilon^2/(2H)$.
\end{proposition}
\begin{proof}
Jensen's inequality gives $\sum_iw_i e^{h_i}\geq
 e^{\sum_iw_i h_i}=1$, proving nonnegativity. Taylor's integral remainder
is the exact identity
\[
 e^x=1+x+x^2\int_0^1(1-t)e^{tx}\,dt.
\]
For $|x|\leq\delta_\varepsilon$, the integral is at most
$e^{\delta_\varepsilon}/2$. Summing the identity with weights $w_i$
cancels the linear term and gives
$\sum_iw_i e^{h_i}\leq1+(e^{\delta_\varepsilon}/2)\sum_iw_i h_i^2$.
The inequality $\log(1+z)\leq z$ for $z\geq0$ follows by integrating
$(1+z)^{-1}\leq1$, and proves the first upper bound. The second uses
$h_i^2\leq\delta_\varepsilon^2$. Monotonicity and Lipschitz continuity of
$P$ give $0\leq P(m+J_\varepsilon)-P(m)\leq LJ_\varepsilon$,
which proves the residual bound for every base $m$.
\end{proof}

For instance, if $h_i=\varepsilon b_i\cdot v+a_i(\varepsilon)$,
$\|v\|\leq M$, $\max_i\|b_i\|\leq B$,
$|a_i(\varepsilon)|\leq C\varepsilon^2$ and the weighted centring
holds exactly, I may take
$\delta_\varepsilon=BM|\varepsilon|+C\varepsilon^2$.
For fixed $B,M,C,L,H$ this yields a uniform $O(\varepsilon^2)$
residual range. Exact centring of the logarithms follows when they are
formed by subtracting their weighted geometric log mean. Centring errors
must otherwise receive their own allowance; they cannot be discarded.

Suppose the implemented output $\widehat X_\varepsilon$ belongs to $[0,1]$,
retained on the $2^{-q}$ grid, and has a proved pointwise normalised
arithmetic error $|\widehat X_\varepsilon(v)-X_\varepsilon(v)|\leq\eta_\varepsilon$ at each supported input. A valid numerical range is
\[
 \rho_\varepsilon=\min\left(1,
 {L e^{\delta_\varepsilon}\delta_\varepsilon^2\over2H}
                  +\eta_\varepsilon\right),
\]
rounded outward and then down to its largest possible grid integer as in
the finite-range certificate. If $\eta_\varepsilon=O(\varepsilon^2)$,
this range remains $O(\varepsilon^2)$ on a bounded support family.
For a nonzero certified integer support upper bound $a_{\max}$, which need not be attained, put
$c=\lceil\log_2 a_{\max}\rceil$ and $R=2^{c-q}$.
The least admissible dyadic range satisfies
$\widehat X_\varepsilon\leq R<2a_{\max}2^{-q}\leq2\rho_\varepsilon$,
with the weak inequality $R\leq2a_{\max}2^{-q}$ also sufficient.
The integer choice is exact: it is the least $c$ with $a_{\max}\leq2^c$.
If $a_{\max}=0$, the mean is zero and needs no measurements.

For an absolute finite-mean tolerance $t_\varepsilon$,
$\E\widehat X_\varepsilon=R p$ is estimated with probability tolerance
$t_\varepsilon/R$ using the range flag. On a sequence with
$t_\varepsilon\geq c_0\varepsilon^2$ and
$R\leq C_0\varepsilon^2$ for positive constants, this probability
tolerance is at least $c_0/C_0$. At fixed confidence within the retained logarithm domain $0<\log(\cdot)<100$, the proved scalar
controller therefore has an $\varepsilon$-independent shot/query cap,
provided this tolerance is in its supported domain; a tolerance at least
half the whole prior gives a zero-shot midpoint instead. To make the controller argument explicit, write $e_R=t_\varepsilon/R$
and let its proved Born prior be $u_R\leq1$. If $u_R\leq2e_R$, the prior
midpoint suffices. Otherwise $u_R>2e_R\geq2c_0/C_0$, so the initial
normalised inverse angle is bounded away from zero and its reciprocal
bootstrap multiplier is bounded. On the supported scalar domain
$10^{-6}\leq e_R<1/2$, the derivative enclosure satisfies $\gamma\leq1$.
The three/five controller's threshold is at most
$\lceil1/(5e_R)\rceil$; its finite shot and geometric-power caps are
therefore bounded at fixed confidence. The full-batch controller's maximum
multiplier is at most $\Pi_+/[4(e_R-10^{-60})]$, where $\Pi_+<4$ is the retained rational upper enclosure of $\pi$; its shot table and admissible
chains are also bounded. These expressions are uniformly finite since
$e_R$ has a positive lower bound in the supported domain. The project choice $\alpha=0.01$ and $e_R\geq10^{-6}$ lies comfortably within the logarithm guard. This is a statement
about the outer estimation factor. The financial evaluator cost, the
numerical register widths and the construction of the support certificate
can still depend on $\varepsilon$.

An unbounded Gaussian law does not satisfy the fixed support hypothesis.
A finite quantile grid does, but its support grows when its address width
grows. Distribution-transfer error and the changing support must then be
charged together with arithmetic precision. Thus the theorem neither
asserts a uniform bound for the continuous Gaussian residual nor proves
constant total pricing cost. It identifies a concrete reason that the
canonical flag's $O(\varepsilon^2)$ variance alone cannot establish a
quantum lower bound or a necessary gain from dyadic value estimation.
